%=====================================================================
\chapter{Instance-Based Learning: k-Nearest Neighbours}\label{ch:knn}
%=====================================================================
\locator{3}{Part III --- Supervised Learning Algorithms}

\begin{outcomes}
By the end of this chapter you will be able to:
\begin{tight}
  \item \textbf{Distinguish} lazy (instance-based) from eager learning.
  \item \textbf{Classify} an example with $k$-nearest neighbours by hand, with
        and without distance weighting.
  \item \textbf{Choose} a distance measure, and \textbf{explain} why features
        must be scaled first.
  \item \textbf{Select} $k$ by cross-validation, relating small and large $k$ to
        variance and bias.
  \item \textbf{Explain} the curse of dimensionality and the computational cost
        of nearest-neighbour search.
\end{tight}
\end{outcomes}

\begin{prereq}
\chref{math} (Euclidean distance), \chref{workflow} (standardisation) and
\chref{generalization} (bias, variance, cross-validation). The HEC outline lists
$k$-NN under unsupervised learning; it is supervised --- it needs labels to vote
with --- and is taught here with the other classifiers.
\end{prereq}

\begin{keyterms}
instance-based learning $\bullet$ lazy vs.\ eager learning $\bullet$
$k$-nearest neighbours $\bullet$ majority vote $\bullet$ distance-weighted vote
$\bullet$ Euclidean, Manhattan, Minkowski and Hamming distance $\bullet$ Voronoi
cell $\bullet$ $k$-NN regression $\bullet$ locally weighted regression $\bullet$
curse of dimensionality $\bullet$ KD-tree $\bullet$ case-based reasoning
\end{keyterms}

\section{Learning by Remembering}

Every learner so far has been \term{eager}: it studies the training data once,
builds an explicit model --- a line, a tree, a table of probabilities --- and then
discards the data. \term{Instance-based} learning is \term{lazy}. It stores the
training examples and does nothing more until asked to classify a new one; only
then does it look for the stored examples most similar to the query and let them
decide.

The idea is old and intuitive: to predict whether a student will pass, find the
students most like them and see how they did. Its strengths follow directly.
There is no training time, the model can be as complex as the data (the decision
boundary follows whatever shape the examples trace), and new examples are added
by simply storing them. Its weaknesses follow too: every prediction requires
searching the stored data, and ``most similar'' is only as good as the distance
measure.

\section{The $k$-Nearest Neighbour Algorithm}

\begin{definitionbox}[$k$-Nearest Neighbours]
To classify a query $\mathbf{x}_q$:
\begin{tightnum}
  \item compute its distance to every stored training example;
  \item take the $k$ examples with the smallest distances;
  \item predict the most common class among them (for regression, the mean of
        their target values).
\end{tightnum}
Its inductive bias is \emph{smoothness}: the classification of an instance is
most similar to the classifications of instances nearby in the feature space.
\end{definitionbox}

\dsfig{%
\begin{tikzpicture}[font=\scriptsize,
  ps/.style={circle, fill=greenhl!70!black, inner sep=1.9pt},
  fl/.style={rectangle, fill=accent, inner sep=2.1pt}]
\draw[-{Stealth[length=2mm]}, gray!60] (0,0) -- (7.4,0) node[below left, font=\tiny] {attendance};
\draw[-{Stealth[length=2mm]}, gray!60] (0,0) -- (0,5.6) node[above right, font=\tiny] {quiz average};
% scale: x = (att-30)/10, y = (quiz-30)/10
\node[ps, label={[font=\tiny]right:A}] at (6.0,5.5) {};
\node[ps, label={[font=\tiny]right:B}] at (5.0,4.0) {};
\node[ps, label={[font=\tiny]above:C}] at (3.0,4.5) {};
\node[fl, label={[font=\tiny]below:D}] at (2.5,1.0) {};
\node[fl, label={[font=\tiny]right:E}] at (4.2,1.6) {};
\node[fl, label={[font=\tiny]left:F}] at (1.0,3.0) {};
\node[star, star points=5, fill=primary, inner sep=1.8pt] (q) at (3.5,3.0) {};
\node[font=\tiny\bfseries, text=primary, anchor=north] at (3.5,2.85) {query};
\draw[primary!60, dashed, line width=0.8pt] (3.5,3.0) circle (1.573);
\draw[secondary!70, dashed, line width=0.8pt] (3.5,3.0) circle (1.82);
\draw[goldhl, dashed, line width=0.8pt] (3.5,3.0) circle (2.52);
\node[font=\tiny, text=primary!80!black, anchor=west] at (7.8,4.6) {$k = 1$: E \quad$\Rightarrow$ Fail};
\node[font=\tiny, text=secondary!70!black, anchor=west] at (7.8,4.0) {$k = 3$: E, C, B \quad$\Rightarrow$ Pass (2--1)};
\node[font=\tiny, text=goldhl!45!black, anchor=west] at (7.8,3.4) {$k = 5$: E, C, B, D, F \quad$\Rightarrow$ Fail (3--2)};
\node[anchor=north west, align=left, text width=5.2cm, font=\tiny, text=gray!55!black] at (7.8,2.9)
     {Circles of increasing radius capture 1, 3 and 5 neighbours. The prediction
      changes twice: $k$ is not a detail but a modelling decision.};
\node[ps] at (8.0,1.1) {}; \node[font=\tiny, anchor=west] at (8.15,1.1) {pass};
\node[fl] at (9.2,1.1) {}; \node[font=\tiny, anchor=west] at (9.35,1.1) {fail};
\end{tikzpicture}}%
{Six students and one query. Which class the query receives depends on how many
neighbours are asked. (Axes scaled so that 10 marks $=$ one unit.)}{knn-votes}

\begin{worked}[Three Values of $k$, Three Answers]
Six students, with attendance and quiz average in percent:
A (90, 85) Pass; B (80, 70) Pass; C (60, 75) Pass; D (55, 40) Fail; E (72, 46) Fail;
F (40, 60) Fail. Classify the query Q (65, 60). Both features are on the same 0--100
scale, so no rescaling is needed.

\smallskip
\textbf{Distances from Q}, sorted:
\begin{tight}
  \item E: $\sqrt{7^2 + 14^2} = \sqrt{245} = 15.65$ \quad (Fail)
  \item C: $\sqrt{5^2 + 15^2} = \sqrt{250} = 15.81$ \quad (Pass)
  \item B: $\sqrt{15^2 + 10^2} = \sqrt{325} = 18.03$ \quad (Pass)
  \item D: $\sqrt{10^2 + 20^2} = \sqrt{500} = 22.36$ \quad (Fail)
  \item F: $\sqrt{25^2 + 0^2} = 25.00$ \quad (Fail)
  \item A: $\sqrt{25^2 + 25^2} = \sqrt{1250} = 35.36$ \quad (Pass)
\end{tight}
\textbf{$k = 1$:} E $\Rightarrow$ \textbf{Fail}. \quad
\textbf{$k = 3$:} E, C, B $\Rightarrow$ \textbf{Pass}, 2 to 1. \quad
\textbf{$k = 5$:} add D, F $\Rightarrow$ \textbf{Fail}, 3 to 2.

\smallskip
\textbf{Distance-weighted, $k = 5$.} Weight each vote by $1/d^2$:
Fail $= \tfrac{1}{245} + \tfrac{1}{500} + \tfrac{1}{625} = 0.00768$;
Pass $= \tfrac{1}{250} + \tfrac{1}{325} = 0.00708$. \textbf{Fail}, narrowly --- E, the
nearest neighbour, counts for more than F, the farthest.
\end{worked}

\section{Distance and Scale}

\begin{definitionbox}[Distance Measures]
For vectors $\mathbf{x}$ and $\mathbf{z}$ with $d$ features:
\begin{tight}
  \item \term{Euclidean}: $\sqrt{\sum_j (x_j - z_j)^2}$ --- the straight-line
        default.
  \item \term{Manhattan}: $\sum_j |x_j - z_j|$ --- less dominated by one large
        difference.
  \item \term{Minkowski}: $\left(\sum_j |x_j - z_j|^p\right)^{1/p}$ --- Manhattan
        at $p = 1$, Euclidean at $p = 2$.
  \item \term{Hamming}: the number of attributes that differ --- for categorical
        data.
  \item \term{Cosine distance}: $1 - \cos\theta$ --- for text and other data where
        direction matters more than length (\chref{math}).
\end{tight}
\end{definitionbox}

\begin{alertbox}[Always Scale Before Measuring Distance]
Add ``logins per week'' (0--40) to the worked example, and a difference of 20 logins
outweighs every difference in marks. The distance is decided by whichever feature
has the largest units. Standardise first (\chref{workflow}). On the breast cancer
data of this chapter's lab, 5-NN makes 7.0\% cross-validated errors on raw
features and 3.3\% after standardising --- the same algorithm, half the error.
\end{alertbox}

A second, subtler problem is \term{irrelevant features}. $k$-NN uses every feature
equally, so twenty features of which two matter give a distance that is mostly
noise. Decision trees choose which attributes to test; $k$-NN cannot, unless the
features are selected or weighted beforehand --- Mitchell's suggestion is to
stretch each axis by a weight chosen by cross-validation.

\section{Choosing $k$}

The decision boundary of 1-NN runs between every pair of differently labelled
neighbours: each training example owns a \term{Voronoi cell}, the region of space
closer to it than to any other example, and the boundary is made of cell edges.
It follows every stray example --- low bias, high variance. Increasing $k$ averages
over more neighbours and smooths the boundary; with $k = n$ it predicts the
majority class everywhere --- high bias, zero variance. $k$ is the regularisation
knob of \chref{generalization}, and it is chosen the same way.

\dsfig{%
\begin{tikzpicture}
\begin{axis}[width=9.8cm, height=5.0cm, axis lines=left, xmode=log, log basis x=10,
  xlabel={$k$ (number of neighbours, log scale)}, ylabel={10-fold CV error},
  tick label style={font=\tiny}, label style={font=\scriptsize},
  xmin=0.8, xmax=190, ymin=0.02, ymax=0.09,
  xtick={1,3,5,11,21,51,101,151}, xticklabels={1,3,5,11,21,51,101,151},
  scaled y ticks=false,
  yticklabel style={/pgf/number format/fixed, /pgf/number format/precision=2}]
\addplot[secondary, line width=1.3pt, mark=*, mark size=1.5pt] coordinates
  {(1,0.0492) (3,0.0352) (5,0.0334) (7,0.0334) (9,0.0334) (11,0.0334) (15,0.0387) (21,0.0440) (31,0.0457) (51,0.0475) (75,0.0563) (101,0.0703) (151,0.0826)};
\node[font=\tiny, accent, anchor=south west, align=left] at (axis cs:1.05,0.05)
     {$k = 1$: follows\\every stray example};
\node[font=\tiny, accent, anchor=south east, align=right] at (axis cs:145,0.083)
     {large $k$: the boundary\\is smoothed away};
\draw[greenhl!70!black, dashed] (axis cs:7,0.02) -- (axis cs:7,0.0334);
\node[font=\tiny, greenhl!45!black, anchor=north west] at (axis cs:7.2,0.031) {best: $k = 5$ to $11$};
\end{axis}
\end{tikzpicture}}%
{Cross-validated error of $k$-NN on the (standardised) breast cancer data. Small $k$
has high variance, large $k$ high bias; the familiar U-shape.}{knn-k}

Two practical rules: for two classes use an odd $k$, so votes cannot tie; and
treat $\sqrt{n}$ as a rough upper starting point, then let cross-validation
decide.

\section{Beyond Classification}

\textbf{$k$-NN regression} predicts the mean (or distance-weighted mean) of the
neighbours' target values. \textbf{Locally weighted regression}, in Mitchell's
treatment, goes further: for each query it fits a small linear model to the
nearby examples only, weighted by their distance, and uses that model's prediction.
The global function is approximated by a patchwork of local ones --- a model that
could never be written down in advance.

\textbf{Case-based reasoning} applies the same lazy idea to rich symbolic
descriptions: a help-desk system that retrieves the most similar past incident and
adapts its solution. Similarity there is measured by matching structure, not by a
formula.

\section{The Curse of Dimensionality}

Nearest-neighbour methods rest on the idea that the nearest examples are
meaningfully closer than the rest. In high dimensions that stops being true.

\dsfig{%
\begin{tikzpicture}
\begin{axis}[width=9.4cm, height=4.8cm, axis lines=left, xmode=log, log basis x=10,
  xlabel={number of dimensions $d$ (log scale)},
  ylabel={nearest / farthest distance},
  tick label style={font=\tiny}, label style={font=\scriptsize},
  xmin=1, xmax=1000, ymin=0, ymax=1.0]
\addplot[accent, line width=1.3pt, mark=*, mark size=1.5pt] coordinates
  {(1,0.001) (2,0.025) (5,0.150) (10,0.293) (20,0.446) (50,0.619) (100,0.712) (200,0.789) (500,0.859) (1000,0.904)};
\addplot[gray!55, dashed] coordinates {(1,1)(1000,1)};
\node[font=\tiny, gray!55!black, anchor=north east, align=right] at (axis cs:900,0.99)
     {1.0: every point equally far away};
\node[font=\tiny, accent, anchor=south east, align=right] at (axis cs:950,0.06)
     {in 1,000 dimensions the nearest of 500\\random points is 90\% as far as the farthest};
\end{axis}
\end{tikzpicture}}%
{The curse of dimensionality. For 500 uniformly random points, the ratio of the
nearest to the farthest distance from a query rises towards 1 as dimensions are
added: ``nearest'' stops meaning anything.}{curse}

\begin{conceptbox}[Why Distances Concentrate]
Each extra dimension adds another independent term to the sum under the square
root. With many terms, the sums for different pairs of points all come out close
to their average --- the same reason an average of many coin tosses is close to
one half. So every point ends up at nearly the same distance from every other, and
the nearest neighbour is barely nearer than a random one. The remedies are fewer
dimensions (feature selection, or PCA in \chref{pcasom}) or a learned notion of
similarity.
\end{conceptbox}

\section{The Cost of Being Lazy}

Training is free; prediction is not. A naive search compares the query with all
$n$ stored examples in $d$ dimensions, $O(nd)$ work per query --- a million
examples means a million distance computations for every prediction.
\term{KD-trees} and ball trees partition space so that most examples can be
ruled out without computing their distance, which works well up to perhaps twenty
dimensions; beyond that, approximate nearest-neighbour indexes trade a little
accuracy for large speed-ups. The whole training set must also be kept, which on
a small device may be impossible.

\begin{pitfall}
\begin{tight}
  \item \textbf{Forgetting to scale.} The feature with the biggest units decides
        everything.
  \item \textbf{Using $k = 1$ on noisy data.} Every mislabelled example claims a
        region of its own.
  \item \textbf{Using an even $k$ for two classes.} Ties must then be broken
        arbitrarily.
  \item \textbf{Keeping irrelevant features.} They add noise to every distance;
        select or weight features first.
  \item \textbf{Ignoring prediction cost.} Fine for thousands of examples, slow for
        millions without an index.
\end{tight}
\end{pitfall}

%---------------------------------------------------------------------
\begin{fourlenses}{Nearest Neighbours}
\lensLA{``Students like you'' is $k$-NN, and it is how many course
recommendation systems work: find the ten most similar past students and suggest
the electives they did well in. The explanation is built in --- here are the
students you resemble --- which advisers find easier to trust than a weight.}

\lensPM{Analogy-based effort estimation is $k$-NN regression: find the three
completed projects most similar to the new one and average their effort. It is
one of the oldest estimation methods in software engineering, and it works
because a team's past projects are the best guide to its next.}

\lensIOT{A machine's current vibration signature is compared with a library of
signatures recorded before known failures; the nearest match names the likely
fault. The library grows every time an engineer labels a new failure, with no
retraining --- the lazy learner's great advantage.}

\lensSEC{Anomaly detection by distance: a login whose nearest neighbours among
the account's own history are all far away is unusual. The curse of dimensionality
is the enemy here --- with hundreds of raw log features, every login looks equally
unusual --- so reduce the features first.}
\end{fourlenses}

%---------------------------------------------------------------------
\begin{lab}[Neighbours, Scaling and the Choice of k]
\begin{lstlisting}[language=Python]
import numpy as np
from sklearn.datasets import load_breast_cancer
from sklearn.model_selection import cross_val_score, GridSearchCV
from sklearn.neighbors import KNeighborsClassifier
from sklearn.pipeline import make_pipeline
from sklearn.preprocessing import StandardScaler

# --- 1. the worked example -----------------------------------------------
X = np.array([[90, 85], [80, 70], [60, 75], [55, 40], [72, 46], [40, 60]])
y = np.array(["Pass", "Pass", "Pass", "Fail", "Fail", "Fail"])
q = np.array([[65, 60]])
for k in (1, 3, 5):
    plain = KNeighborsClassifier(k).fit(X, y).predict(q)[0]
    weighted = KNeighborsClassifier(k, weights="distance").fit(X, y)
    print(f"k={k}: vote {plain}, distance-weighted {weighted.predict(q)[0]}")

# --- 2. scaling matters ----------------------------------------------------
Xb, yb = load_breast_cancer(return_X_y=True)
raw = cross_val_score(KNeighborsClassifier(5), Xb, yb, cv=10).mean()
scaled = cross_val_score(make_pipeline(StandardScaler(),
                                       KNeighborsClassifier(5)),
                         Xb, yb, cv=10).mean()
print(f"5-NN accuracy: raw {raw:.3f}, standardised {scaled:.3f}")

# --- 3. choose k by cross-validation --------------------------------------
pipe = make_pipeline(StandardScaler(), KNeighborsClassifier())
grid = GridSearchCV(pipe, {"kneighborsclassifier__n_neighbors":
                           [1, 3, 5, 7, 11, 21, 51, 101]}, cv=10)
grid.fit(Xb, yb)
for k, s in zip([1, 3, 5, 7, 11, 21, 51, 101],
                grid.cv_results_["mean_test_score"]):
    print(f"  k={k:3d}: CV accuracy {s:.3f}")

# --- 4. the curse of dimensionality ----------------------------------------
rng = np.random.default_rng(0)
for dim in (2, 10, 100, 1000):
    P, Q = rng.uniform(size=(500, dim)), rng.uniform(size=(50, dim))
    D = np.sqrt(((Q[:, None, :] - P[None, :, :]) ** 2).sum(-1))
    ratio = (D.min(1) / D.max(1)).mean()
    print(f"d={dim:4d}: nearest/farthest = {ratio:.2f}")
\end{lstlisting}
scikit-learn's \texttt{weights="distance"} uses $1/d$ rather than the $1/d^2$ of
the worked example; here both give the same answers. Add ten columns of random
noise to the breast cancer features and rerun part 3: what happens to the best
accuracy?
\end{lab}

%---------------------------------------------------------------------
\begin{checkpoint}
\begin{tightnum}
  \item In the worked example, what does 3-NN predict for the query (50, 55)?
  \item Why is $k$-NN called a lazy learner, and what does laziness cost at
        prediction time?
  \item As $k$ grows from 1 to $n$, what happens to bias and to variance?
\end{tightnum}
\end{checkpoint}

%---------------------------------------------------------------------
\begin{chaptersummary}
\begin{tight}
  \item $k$-NN stores the training data and classifies a query by a vote of its
        $k$ nearest neighbours (the mean, for regression). Its bias is
        smoothness: nearby examples share labels.
  \item The prediction can change with $k$; distance weighting lets closer
        neighbours count more.
  \item Scale features before computing distances, and remove irrelevant ones;
        choose the distance to suit the data.
  \item Small $k$ means high variance, large $k$ high bias; choose $k$ by
        cross-validation, odd for two classes.
  \item In high dimensions all distances become similar, and nearest neighbours
        stop being meaningful.
  \item Training is free but each prediction searches the data; indexes such as
        KD-trees reduce the cost in low dimensions.
\end{tight}
\end{chaptersummary}

\begin{reviewq}
  \item \textbf{(MCQ)} $k$-NN is: (a)~an eager learner (b)~a lazy learner
        (c)~an unsupervised method (d)~a linear model.
  \item \textbf{(MCQ)} Increasing $k$ in $k$-NN generally: (a)~raises variance
        (b)~lowers bias (c)~smooths the decision boundary (d)~makes training
        slower.
  \item \textbf{(MCQ)} The Manhattan distance between $(1, 2)$ and $(4, 6)$ is:
        (a)~5 (b)~7 (c)~25 (d)~12.
  \item \textbf{(Short)} Explain the curse of dimensionality and its effect on
        $k$-NN.
  \item \textbf{(Short)} Why does $k$-NN need feature scaling while a decision tree
        does not?
  \item \textbf{(Short)} Describe locally weighted regression in your own words.
  \item \textbf{(Applied)} Add a seventh student G (68, 58, Pass) to the worked
        example. Recompute the 1-, 3- and 5-NN predictions for Q.
  \item \textbf{(Applied)} A dataset has features ``age'' (18--60) and ``annual
        income'' (100{,}000--2{,}000{,}000 rupees). Compute the Euclidean distance
        between (25, 400{,}000) and (55, 410{,}000) before and after min--max
        scaling, and comment.
\end{reviewq}
